\chapter{Introducción}
\label{ch:intro}

\section{Dominio del problema}

GraphHunter es una herramienta interactiva de \emph{threat hunting}
construida alrededor de una idea: la mayor parte del comportamiento
de un atacante es naturalmente un camino en un grafo dirigido
temporal, así que dejamos al analista \emph{pedir caminos} en vez
de escribir SQL por tabla. Una hipótesis se expresa en un pequeño
DSL de cadena de flechas --- por ejemplo,

\begin{center}
\code{User -[Auth \{status="Failure"\}]-> IP HAVING count(distinct IP)
>= 3 WITHIN 24h}
\end{center}

y el motor devuelve los caminos del grafo ingerido que la
satisfacen, rankeados por un puntaje de anomalía compuesto. El
grafo se construye a partir de fuentes de logs heterogéneas
(Microsoft Sentinel, Defender XDR, Fortinet, Sysmon, IIS, CSV/JSON
genéricos). Las cargas típicas van desde unos pocos miles de
aristas (un ticket de incidente) hasta cientos de millones (dump
EVTX de varias semanas de una empresa mediana).

\section{Vista general del sistema}

El código está organizado como un monorepo Rust con cinco
componentes principales, dispuestos en capas como en la
Figura~\ref{fig:crate-graph}.

\begin{figure}[h]
\centering
\begin{tikzpicture}[
  node distance=4mm and 10mm,
  crate/.style={draw, rounded corners, minimum width=3.2cm,
                minimum height=0.85cm, align=center, font=\small},
  layer/.style={draw, dashed, rounded corners, inner sep=4mm},
]

  \node[crate] (core) {\textbf{graph\_hunter\_core}\\ motor};
  \node[crate, right=of core] (api) {\textbf{graph\_hunter\_api}\\ fachada canónica};

  \node[crate, above right=of api] (tauri) {Tauri desktop\\ (\texttt{app/src-tauri})};
  \node[crate, right=of api]       (http)  {Servidor HTTP\\ (\texttt{http\_api.rs})};
  \node[crate, below right=of api] (mcp)   {Servidor MCP\\ (\texttt{graph-hunter-mcp})};

  \node[crate, below=of core] (libgm) {\textbf{libgraphmatch}\\ SIMD AVX2 en C++};

  \draw[-{Stealth}] (core) -- (api);
  \draw[-{Stealth}] (api) -- (tauri);
  \draw[-{Stealth}] (api) -- (http);
  \draw[-{Stealth}] (api) -- (mcp);
  \draw[-{Stealth}] (libgm) -- (core) node[midway,right,font=\tiny]{FFI};
  \draw[-{Stealth}, dashed] (mcp) to[bend right=25] node[above,font=\tiny]{HTTP} (http);

\end{tikzpicture}
\caption{Grafo de crates. El motor (\texttt{graph\_hunter\_core})
lo consume una fachada canónica (\texttt{graph\_hunter\_api});
tres capas de transporte delegan en esa fachada. El servidor MCP
llega al mismo motor invocando el servidor HTTP sobre la interfaz
de loopback, en vez de enlazarse directamente a la fachada.}
\label{fig:crate-graph}
\end{figure}

El motor es una biblioteca Rust pura: sin dependencias de
framework, sin estado global salvo un puñado de catálogos cacheados
vía \texttt{OnceLock}. La crate fachada (\code{graph\_hunter\_api})
envuelve al motor con gestión de sesión, DTOs y un trait
\code{EventEmitter}, de modo que las tres capas de transporte
compartan una única implementación sin importar los tipos de bajo
nivel. El backend SIMD en C++ (\code{libgraphmatch}) se enlaza por
FFI cuando el feature \code{simd} está activado; una alternativa
pura en Rust (\code{simd\_rust}) ya está en camino de reemplazarlo
(Capítulo~\ref{ch:simd}).

\section{Qué cubre esta referencia}

El documento recorre el motor de abajo hacia arriba:

\begin{itemize}
    \item \textbf{Parte II} --- las estructuras de datos que el
        motor mantiene en memoria: relaciones compactas, el
        internador de strings bidireccional, el grafo en memoria,
        el edge store con capacidad de \emph{spill}, el almacén
        externo de metadatos y el AST de hipótesis que produce el
        parser del DSL.
    \item \textbf{Parte III} --- los algoritmos principales:
        parseo del DSL, matcheo de patrones temporales (DFS con
        podado, variante \emph{smart} con cota de anomalía),
        aceleración SIMD, \emph{scoring} de anomalía, analíticos
        de centralidad (betweenness de Brandes, PageRank
        temporal), deduplicación \& paginación de caminos, y
        \emph{diff} de hunts.
    \item \textbf{Parte IV} --- el pipeline de ingesta en
        \emph{streaming}: trait \code{LogSource}, parsers,
        escritor con canales acotados, estampado MVCC a nivel de
        arista y persistencia de desborde.
    \item \textbf{Parte V} --- los puntos de extensión que
        introdujo el refactor P2: \code{GraphRead}/\code{GraphWrite},
        \code{ScoreComponent}/\code{Scorer},
        \code{StringInternerBackend}, \code{IngestAdapter} y el
        cargador de catálogo YAML.
    \item \textbf{Parte VI} --- la capa de transporte: la API
        canónica, los DTOs, \code{EventEmitter}, los adaptadores
        Tauri/HTTP/MCP, el arnés de paridad y la persistencia de
        sesión.
\end{itemize}

Cierran el documento cuatro apéndices: un resumen de complejidad
de una página (Ap.~\ref{app:complexity}), la referencia del
catálogo ATT\&CK (Ap.~\ref{app:catalog}), una tabla de
benchmarks (Ap.~\ref{app:bench}) y un índice de código que mapea
cada algoritmo a su \code{archivo:línea} (Ap.~\ref{app:source}).

\begin{inthecode}
Las raíces de crates son \filepath{graph\_hunter\_core/Cargo.toml},
\filepath{graph\_hunter\_api/Cargo.toml},
\filepath{app/src-tauri/Cargo.toml},
\filepath{graph-hunter-mcp/package.json} y
\filepath{libgraphmatch/CMakeLists.txt}.
\end{inthecode}
